\documentclass[11pt]{article}
\usepackage{amsmath,amssymb,amsthm,booktabs}
\usepackage[margin=1in]{geometry}

\newcommand{\MGamma}{\mathcal M_\Gamma}
\newcommand{\Uinf}{\mathcal U_\infty}
\newcommand{\Kamb}{K_a^{\Gamma,\mathrm{amb}}}

\title{Step 178: Derivation Attempt for \(\kappa_{a,w}\)}
\author{Six Birds Foundations III}
\date{2026-05-15}

\begin{document}
\maketitle

\section{Verdict}

\[
\boxed{\texttt{kappa\_verdict = V\_kappa\_classical\_theorem\_needed}.}
\]

The target vector is
\[
\kappa_{a,w}(\tau)
=
\left(T_a^*K_a^\Gamma(\cdot,w)\right)(1/2+i\tau).
\]
For \(a=1/2\) and \(w=\rho_1=1/2+14.134725141734693i\), the inherited
records reduce the problem to one specific missing classical theorem:

\[
\boxed{
\text{Burnol's explicit Bessel/Hankel resolvent formula for }
P_{\mathcal L_a^\Gamma}
\text{, equivalently }K_a^\Gamma.
}
\]

\section{T1. Inherited Records}

Step 145 records the Burnol carrier:
\[
L_a\subset L^2(0,\infty;dt)
\]
is Burnol's extended Sonine space, consisting of functions constant on
\((0,a)\) whose cosine transform is also constant on \((0,a)\).

Step 153 records
\[
T_a:=\MGamma J_a\Uinf^{-1},
\]
and
\[
\Uinf(J_a^*Y^a_{w,k})
=
T_a^*\partial_{\overline w}^{\,k}K_a^\Gamma(\cdot,w).
\]
Thus, for \(k=0\),
\[
\kappa_{a,w}=T_a^*K_a^\Gamma(\cdot,w).
\]

Step 153 also records the completed ambient Hardy kernel
\[
\Kamb(s,w)
=
A_\infty(s)\overline{A_\infty(w)}
\frac{s}{s-1}\overline{\frac{w}{w-1}}
\frac{a^{s+\overline w-1}}{s+\overline w-1}.
\]
The exact Burnol/Sonine kernel is only given as
\[
\boxed{
K_a^\Gamma(\cdot,w)
=
P_{\mathcal L_a^\Gamma}\Kamb(\cdot,w).
}
\]
Step 153 explicitly states that this projection is load-bearing; replacing it
by the ambient Hardy shadow is not lawful.

Step 152 records the adjoint transport context
\[
\mathfrak B_{\ell,a}^*Y^a_{\rho,k}
=
(I-P_\infty)\tau_{-\ell}P_\infty J_a^*Y^a_{\rho,k}.
\]
Step 113 records Burnol's complement theorem for \(a<1\):
\[
Y_a^\perp=P_a.
\]

A repository search found no root-level \texttt{adequacy.tex} or
\texttt{needles.tex} files.

\section{T2. Derivation Attempts}

\subsection*{Attack A: Closed-form projected-kernel route}

Starting from Step 153,
\[
\kappa_{a,w}
=
T_a^*K_a^\Gamma(\cdot,w)
=
T_a^*P_{\mathcal L_a^\Gamma}\Kamb(\cdot,w).
\]
The ambient kernel \(\Kamb\) is explicit.  Therefore a closed form for
\(\kappa\) would follow from an explicit formula for
\[
P_{\mathcal L_a^\Gamma}\Kamb(\cdot,w).
\]
This is precisely the projected Sonine reproducing kernel.  No inherited
record supplies this projection.

\subsection*{Attack B: Transport-kernel route}

If the adjoint transport had an integral kernel \(T_a^*(\tau,z)\), then
\[
\boxed{
\kappa_{a,w}(\tau)
=
\int \overline{T_a(z,\tau)}\,K_a^\Gamma(z,w)\,d\mu_a(z).
}
\]
Substituting the projected kernel gives
\[
\kappa_{a,w}(\tau)
=
\int \overline{T_a(z,\tau)}
\left(P_{\mathcal L_a^\Gamma}\Kamb(\cdot,w)\right)(z)
d\mu_a(z).
\]
This is an operational identity, but it is not numerically evaluable because
neither the \(T_a\) integral kernel nor the projected kernel action is recorded
in inherited artifacts.

\subsection*{Attack C: Spectral route}

If a Bessel/Hankel spectral theorem supplied an orthonormal basis
\(\{u_r\}\) of \(\mathcal L_a^\Gamma\), then
\[
P_{\mathcal L_a^\Gamma}
=
\sum_r |u_r\rangle\langle u_r|
\]
and
\[
\boxed{
\kappa_{a,w}(\tau)
=
\sum_r
\langle \Kamb(\cdot,w),u_r\rangle
(T_a^*u_r)(1/2+i\tau).
}
\]
This is the desired spectral expansion.  It requires the same missing
classical input: the Bessel/Hankel spectral data for the projected Sonine
space \(L_a\).

\section{Specific Classical Theorem Needed}

The needed theorem is not generic ``operator theory.'' It is:

\begin{quote}
Burnol's explicit Bessel/Hankel resolvent formula for the orthogonal
projection \(P_{\mathcal L_a^\Gamma}\), equivalently the reproducing kernel
\(K_a^\Gamma=P_{\mathcal L_a^\Gamma}K_a^{\Gamma,\mathrm{amb}}\), for the
extended Sonine space \(L_a\) at \(a=1/2\).
\end{quote}

This is the theorem that would turn the ambient Hardy kernel into the legal
projected Burnol/Sonine kernel and then allow \(T_a^*\) to be applied.

\section{T3--T4. Specialization and Samples}

For
\[
a=\frac12,\qquad w=\rho_1=\frac12+14.134725141734693i,
\]
the formal expression is
\[
\kappa_{1/2,\rho_1}(\tau)
=
T_{1/2}^*
P_{\mathcal L_{1/2}^\Gamma}
K_{1/2}^{\Gamma,\mathrm{amb}}(\cdot,\rho_1)(1/2+i\tau).
\]
No sample values are lawfully computed.  The requested sample set
\[
\tau\in\{-14.134725141734693,-5,-1,0,1,5,14.134725141734693\}
\]
is recorded in \texttt{kappa\_sample\_values\_step178.csv} with status
\[
\texttt{not\_computed\_classical\_projection\_theorem\_needed}.
\]

\section{T6. \(c_{11}(\log2)\)}

Since \(\kappa_{1/2,\rho_1}\) is not derived, the Step 177 formula
\[
c_{11}(\ell)
=
\frac{\langle \kappa_1-P\kappa_1,
M_{m_\ell}P\kappa_1\rangle}{\|\kappa_1\|^2}
\]
cannot be evaluated.  The row for \(c_{11}(\log2)\) is recorded in
\texttt{c\_ij\_attempt\_step178.csv} with status
\[
\texttt{not\_attempted\_kappa\_unavailable}.
\]

\section{Branch B Implication}

Step 178 advances Branch B by replacing the Step 177 transport ambiguity with
a named classical-content interface.  The exact item needed is the projected
Sonine reproducing kernel / Bessel-Hankel resolvent theorem.  Until it is
supplied, SL164.1 and the finite matrix entries remain blocked.

\section{Nonclaim Boundary}

This step does not prove RH, does not close \(\Xi_{\mathrm{BC}}\), does not
close Branch B, does not compute \(\kappa\), and does not compute
\(c_{11}(\log2)\).  It does not substitute the ambient Hardy kernel for the
projected Sonine kernel.

\end{document}
